%% kapa: exam=1 round=19 year=2008 period=1 form=A subject=accounting subject_ko=회계학 seq=20 number=100
%% source: Q-Net 감정평가사 기출문제 (cst00309) artlSeq=1001962 file=extracted/1008/1001962/2008년 제19회 감정평가사 1차시험 기출문제/1교시.hwp
%% conversion: kordoc 4.24.51 -> tools/gen1.py (auto); review=proofread-sonnet; flags=none
\begin{question}{100}
(주)광림중공업은 선박용 엔진을 제작하여 조선회사에 장기할부판매를 하고 있다. (주)광림중공업은 20X1년 7월 1일 (주)동원조선에 엔진 1대를 ￦1,200,000에 할부로 판매하고, 판매 시 현금 ￦200,000을 받고 잔액은 향후 2년간에 걸쳐 매년 7월 1일에 ￦500,000씩 받기로 하였다. 장기매출채권에 대한 할인율은 연 10\%이고, 장기매출채권의 회수가능성이 매우 높다. (주)광림중공업은 중소기업이 아니며, 장기매출채권의 명목가액과 장부금액 간의 차이가 중요하다고 가정할 경우 (주)광림중공업이 판매시점에 계상할 장기매출채권의 장부금액은 얼마인가? (단, 할인율 10\%에 대한 기간 2년의 단일금액현가계수와 연금현가계수는 각각 0.8264와 1.7355이다.)
\begin{choices}
  \item ￦867,750
  \item ￦1,000,000
  \item ￦1,200,000
  \item ￦1,735,500
  \item ￦2,082,600
\end{choices}
\end{question}
\answer{1}
